\begin{examplebox}
  \textbf{\textcolor{CExample}{例 1｜知道方向，直接写新函数（基础）}}

  \textbf{题目：}把 $y=x^2$ 的图像向左平移 $3$ 个单位，写出新解析式。

  \textbf{最快做法：}左移 $3$ 格，就把原式的 $x$ 换为 $x+3$：
  \[
    \boxed{g(x)=(x+3)^2}.
  \]
  \textbf{为什么对：}原来顶点在 $(0,0)$，现在应该在 $(-3,0)$；代入 $x=-3$，恰好有 $g(-3)=0$。
  \begin{center}
  \begin{tikzpicture}[x=0.85cm,y=0.7cm,>=stealth,font=\small]
    \draw[->,gray] (-4.7,0)--(1.8,0) node[right] {$x$};
    \draw[->,gray] (0,-0.25)--(0,2.75) node[above] {$y$};
    \draw[very thick,blue!75!black,domain=-1.55:1.55,samples=55] plot (\x,{\x*\x});
    \draw[very thick,teal!75!black,domain=-4.55:-1.45,samples=55] plot (\x,{(\x+3)*(\x+3)});
    \fill[teal!75!black] (-3,0) circle (2pt);
    \fill[blue!75!black] (0,0) circle (2pt);
    \node[below] at (-3,0) {$-3$};
    \node[below right] at (0,0) {$0$};
    \draw[->,teal!75!black,very thick] (-0.25,-0.60)--(-2.75,-0.60)
      node[midway,below] {左移 $3$ 格};
    \node[teal!75!black,above] at (-3.5,2.2) {$g(x)=(x+3)^2$};
    \node[blue!75!black,above] at (0.65,2.2) {$f(x)=x^2$};
  \end{tikzpicture}

  \end{center}

  \medskip
  \textbf{\textcolor{CExample}{例 2｜知道新旧函数，反求距离（常考）}}

  \textbf{题目：}将 $f(x)=3x-2$ 的图像向右平移 $a$ 个单位后，得到 $g(x)=3x-8$。求 $a$。

  \textbf{最快做法：}先追踪高度都是 $-2$ 的点。
  原图像过 $(0,-2)$，新图像过 $(2,-2)$，横坐标增大 $2$，因此\textbf{$a=2$}。

  \textbf{代数核对：}
  \[
    g(x)=f(x-a)=3(x-a)-2=3x-3a-2.
  \]
  与 $3x-8$ 比较，$-3a-2=-8$，得\textbf{$\boxed{a=2}$}。
  \begin{center}
  \begin{tikzpicture}[x=0.98cm,y=0.44cm,>=stealth,font=\small]
    \draw[->,gray] (-0.45,0)--(4.2,0) node[right] {$x$};
    \draw[->,gray] (0,-3.1)--(0,4.1) node[above] {$y$};
    \draw[very thick,blue!75!black,domain=-0.25:1.95,samples=2] plot (\x,{3*\x-2});
    \draw[very thick,teal!75!black,domain=1.75:3.95,samples=2] plot (\x,{3*\x-8});
    \draw[densely dashed,gray] (0,-2)--(2,-2);
    \fill[blue!75!black] (0,-2) circle (2.2pt);
    \fill[teal!75!black] (2,-2) circle (2.2pt);
    \node[below left,blue!75!black] at (0,-2) {$(0,-2)$};
    \node[below right,teal!75!black] at (2,-2) {$(2,-2)$};
    \draw[->,teal!75!black,thick] (0.32,-1.5)--(1.72,-1.5)
      node[midway,above] {右移 $2$};
    \node[blue!75!black,above] at (1.0,4.6) {$f(x)$};
    \node[teal!75!black,above] at (3.2,4.6) {$g(x)$};
  \end{tikzpicture}

  \end{center}

  \medskip
  \textbf{\textcolor{CExample}{例 3｜利用对称轴，少算很多步（提速）}}

  \textbf{题目：}$f(x)=x^2-2x$ 向右平移 $t$ 个单位后，新图像的对称轴是 $x=3$，求 $t$。

  \textbf{最快做法：}先配方：
  \[
    f(x)=(x-1)^2-1,
  \]
  原对称轴是 $x=1$。现在轴到了 $x=3$，所以向右移了 $3-1=2$ 格，\textbf{$\boxed{t=2}$}。

  \textbf{为什么比展开快：}整条图像平移时，顶点和对称轴一起移动；不必先把新解析式展开。
  \begin{center}
  \begin{tikzpicture}[x=0.9cm,y=0.6cm,>=stealth,font=\small]
    \draw[->,gray] (-0.55,0)--(5.4,0) node[right] {$x$};
    \draw[->,gray] (0,-1.55)--(0,2.55) node[above] {$y$};
    \draw[thick,blue!75!black,domain=-0.55:2.75,samples=55] plot (\x,{(\x-1)*(\x-1)-1});
    \draw[thick,teal!75!black,domain=1.25:4.75,samples=55] plot (\x,{(\x-3)*(\x-3)-1});
    \draw[densely dashed,blue!75!black] (1,-1.35)--(1,2.25);
    \draw[densely dashed,teal!75!black] (3,-1.35)--(3,2.25);
    \fill[blue!75!black] (1,-1) circle (2pt);
    \fill[teal!75!black] (3,-1) circle (2pt);
    \node[below left,blue!75!black] at (1,-1) {$(1,-1)$};
    \node[below right,teal!75!black] at (3,-1) {$(3,-1)$};
    \node[above,blue!75!black] at (1,2.25) {$x=1$};
    \node[above,teal!75!black] at (3,2.25) {$x=3$};
    \draw[->,teal!75!black,very thick] (1.4,-1.65)--(2.6,-1.65)
      node[midway,below] {右移 $2$};
  \end{tikzpicture}

  \end{center}

  \textbf{\textcolor{CExample}{考场判断顺序：}}先看图像上的\textbf{顶点、对称轴或对应点}；看不出来时，再用 $f(x+a)$ 或 $f(x-a)$ 代入。
\end{examplebox}
